\documentclass[12pt]{article}

% Idioma y codificación
\usepackage[spanish, es-tabla]{babel}       %es-tabla para que se titule "Tabla"
\usepackage[utf8]{inputenc}

% Márgenes
\usepackage[a4paper,top=3cm,bottom=2.5cm,left=3cm,right=3cm]{geometry}

% Comentarios de bloque
\usepackage{verbatim}

% Paquetes de links
\usepackage[hidelinks]{hyperref}    % Permite enlaces
\usepackage{url}                    % redirecciona a la web

% Más opciones para enumeraciones
\usepackage{enumitem}

% Personalizar la portada
\usepackage{titling}


% Paquetes de tablas
\usepackage{multirow}


%------------------------------------------------------------------------

%Paquetes de figuras
\usepackage{caption}
\usepackage{subcaption} % Figuras al lado de otras
\usepackage{float}      % Poner figuras en el sitio indicado H.


% Paquetes de imágenes
\usepackage{graphicx}       % Paquete para añadir imágenes
\usepackage{transparent}    % Para manejar la opacidad de las figuras

% Paquete para usar colores
\usepackage[dvipsnames]{xcolor}
\usepackage{pagecolor}      % Para cambiar el color de la página

% Habilita tamaños de fuente mayores
\usepackage{fix-cm}


%------------------------------------------------------------------------

% Paquetes de matemáticas
\usepackage{mathtools, amsfonts, amssymb, mathrsfs}
\usepackage[makeroom]{cancel}     % Simplificar tachando
\usepackage{polynom}    % Divisiones y Ruffini
\usepackage{units} % Para poner fracciones diagonales con \nicefrac

\usepackage{pgfplots}   %Representar funciones
\pgfplotsset{compat=1.18}  % Versión 1.18

\usepackage{tikz-cd}    % Para usar diagramas de composiciones
\usetikzlibrary{calc}   % Para usar cálculo de coordenadas en tikz

%Definición de teoremas, etc.
\usepackage{amsthm}
%\swapnumbers   % Intercambia la posición del texto y de la numeración

\theoremstyle{plain}

\makeatletter
\@ifclassloaded{article}{
  \newtheorem{teo}{Teorema}[section]
}{
  \newtheorem{teo}{Teorema}[chapter]  % Se resetea en cada chapter
}
\makeatother

\newtheorem{coro}{Corolario}[teo]           % Se resetea en cada teorema
\newtheorem{prop}[teo]{Proposición}         % Usa el mismo contador que teorema
\newtheorem{lema}[teo]{Lema}                % Usa el mismo contador que teorema

\theoremstyle{remark}
\newtheorem*{observacion}{Observación}

\theoremstyle{definition}

\makeatletter
\@ifclassloaded{article}{
  \newtheorem{definicion}{Definición} [section]     % Se resetea en cada chapter
}{
  \newtheorem{definicion}{Definición} [chapter]     % Se resetea en cada chapter
}
\makeatother

\newtheorem*{notacion}{Notación}
\newtheorem*{ejemplo}{Ejemplo}
\newtheorem*{ejercicio*}{Ejercicio}             % No numerado
\newtheorem{ejercicio}{Ejercicio} [section]     % Se resetea en cada section


% Modificar el formato de la numeración del teorema "ejercicio"
\renewcommand{\theejercicio}{%
  \ifnum\value{section}=0 % Si no se ha iniciado ninguna sección
    \arabic{ejercicio}% Solo mostrar el número de ejercicio
  \else
    \thesection.\arabic{ejercicio}% Mostrar número de sección y número de ejercicio
  \fi
}


% \renewcommand\qedsymbol{$\blacksquare$}         % Cambiar símbolo QED
%------------------------------------------------------------------------

% Paquetes para encabezados
\usepackage{fancyhdr}
\pagestyle{fancy}
\fancyhf{}

\newcommand{\helv}{ % Modificación tamaño de letra
\fontfamily{}\fontsize{12}{12}\selectfont}
\setlength{\headheight}{15pt} % Amplía el tamaño del índice


%\usepackage{lastpage}   % Referenciar última pag   \pageref{LastPage}
\fancyfoot[C]{\thepage}

%------------------------------------------------------------------------

% Conseguir que no ponga "Capítulo 1". Sino solo "1."
\makeatletter
\@ifclassloaded{book}{
  \renewcommand{\chaptermark}[1]{\markboth{\thechapter.\ #1}{}} % En el encabezado
    
  \renewcommand{\@makechapterhead}[1]{%
  \vspace*{50\p@}%
  {\parindent \z@ \raggedright \normalfont
    \ifnum \c@secnumdepth >\m@ne
      \huge\bfseries \thechapter.\hspace{1em}\ignorespaces
    \fi
    \interlinepenalty\@M
    \Huge \bfseries #1\par\nobreak
    \vskip 40\p@
  }}
}
\makeatother

%------------------------------------------------------------------------
% Paquetes de cógido
\usepackage{minted}
\renewcommand\listingscaption{Código fuente}

\usepackage{fancyvrb}
% Personaliza el tamaño de los números de línea
\renewcommand{\theFancyVerbLine}{\small\arabic{FancyVerbLine}}

% Estilo para C++
\newminted{cpp}{
    frame=lines,
    framesep=2mm,
    baselinestretch=1.2,
    linenos,
    escapeinside=||
}



\usepackage{listings} % Para incluir código desde un archivo

\renewcommand\lstlistingname{Código Fuente}
\renewcommand\lstlistlistingname{Índice de Códigos Fuente}

% Definir colores
\definecolor{vscodepurple}{rgb}{0.5,0,0.5}
\definecolor{vscodeblue}{rgb}{0,0,0.8}
\definecolor{vscodegreen}{rgb}{0,0.5,0}
\definecolor{vscodegray}{rgb}{0.5,0.5,0.5}
\definecolor{vscodebackground}{rgb}{0.97,0.97,0.97}
\definecolor{vscodelightgray}{rgb}{0.9,0.9,0.9}

% Configuración para el estilo de C similar a VSCode
\lstdefinestyle{vscode_C}{
  backgroundcolor=\color{vscodebackground},
  commentstyle=\color{vscodegreen},
  keywordstyle=\color{vscodeblue},
  numberstyle=\tiny\color{vscodegray},
  stringstyle=\color{vscodepurple},
  basicstyle=\scriptsize\ttfamily,
  breakatwhitespace=false,
  breaklines=true,
  captionpos=b,
  keepspaces=true,
  numbers=left,
  numbersep=5pt,
  showspaces=false,
  showstringspaces=false,
  showtabs=false,
  tabsize=2,
  frame=tb,
  framerule=0pt,
  aboveskip=10pt,
  belowskip=10pt,
  xleftmargin=10pt,
  xrightmargin=10pt,
  framexleftmargin=10pt,
  framexrightmargin=10pt,
  framesep=0pt,
  rulecolor=\color{vscodelightgray},
  backgroundcolor=\color{vscodebackground},
}

%------------------------------------------------------------------------

% Comandos definidos
\newcommand{\bb}[1]{\mathbb{#1}}
\newcommand{\cc}[1]{\mathcal{#1}}

% I prefer the slanted \leq
\let\oldleq\leq % save them in case they're every wanted
\let\oldgeq\geq
\renewcommand{\leq}{\leqslant}
\renewcommand{\geq}{\geqslant}

% Si y solo si
\newcommand{\sii}{\iff}

% Letras griegas
\newcommand{\eps}{\epsilon}
\newcommand{\veps}{\varepsilon}
\newcommand{\lm}{\lambda}

\newcommand{\ol}{\overline}
\newcommand{\ul}{\underline}
\newcommand{\wt}{\widetilde}
\newcommand{\wh}{\widehat}

\let\oldvec\vec
\renewcommand{\vec}{\overrightarrow}

% Derivadas parciales
\newcommand{\del}[2]{\frac{\partial #1}{\partial #2}}
\newcommand{\Del}[3]{\frac{\partial^{#1} #2}{\partial^{#1} #3}}
\newcommand{\deld}[2]{\dfrac{\partial #1}{\partial #2}}
\newcommand{\Deld}[3]{\dfrac{\partial^{#1} #2}{\partial^{#1} #3}}


\newcommand{\AstIg}{\stackrel{(\ast)}{=}}
\newcommand{\Hop}{\stackrel{L'H\hat{o}pital}{=}}

\newcommand{\red}[1]{{\color{red}#1}} % Para integrales, destacar los cambios.

% Método de integración
\newcommand{\MetInt}[2]{
    \left[\begin{array}{c}
        #1 \\ #2
    \end{array}\right]
}

% Declarar aplicaciones
% 1. Nombre aplicación
% 2. Dominio
% 3. Codominio
% 4. Variable
% 5. Imagen de la variable
\newcommand{\Func}[5]{
    \begin{equation*}
        \begin{array}{rrll}
            #1:& #2 & \longrightarrow & #3\\
               & #4 & \longmapsto & #5
        \end{array}
    \end{equation*}
}

%------------------------------------------------------------------------


\begin{document}

    % 1. Foto de fondo
    % 2. Título
    % 3. Encabezado Izquierdo
    % 4. Color de fondo
    % 5. Coord x del titulo
    % 6. Coord y del titulo
    % 7. Fecha

    
    % 1. Foto de fondo
% 2. Título
% 3. Encabezado Izquierdo
% 4. Color de fondo
% 5. Coord x del titulo
% 6. Coord y del titulo
% 7. Fecha

\newcommand{\portada}[7]{

    \portadaBase{#1}{#2}{#3}{#4}{#5}{#6}{#7}
    \portadaBook{#1}{#2}{#3}{#4}{#5}{#6}{#7}
}

\newcommand{\portadaExamen}[7]{

    \portadaBase{#1}{#2}{#3}{#4}{#5}{#6}{#7}
    \portadaArticle{#1}{#2}{#3}{#4}{#5}{#6}{#7}
}




\newcommand{\portadaBase}[7]{

    % Tiene la portada principal y la licencia Creative Commons
    
    % 1. Foto de fondo
    % 2. Título
    % 3. Encabezado Izquierdo
    % 4. Color de fondo
    % 5. Coord x del titulo
    % 6. Coord y del titulo
    % 7. Fecha
    
    
    \thispagestyle{empty}               % Sin encabezado ni pie de página
    \newgeometry{margin=0cm}        % Márgenes nulos para la primera página
    
    
    % Encabezado
    \fancyhead[L]{\helv #3}
    \fancyhead[R]{\helv \nouppercase{\leftmark}}
    
    
    \pagecolor{#4}        % Color de fondo para la portada
    
    \begin{figure}[p]
        \centering
        \transparent{0.3}           % Opacidad del 30% para la imagen
        
        \includegraphics[width=\paperwidth, keepaspectratio]{assets/#1}
    
        \begin{tikzpicture}[remember picture, overlay]
            \node[anchor=north west, text=white, opacity=1, font=\fontsize{60}{90}\selectfont\bfseries\sffamily, align=left] at (#5, #6) {#2};
            
            \node[anchor=south east, text=white, opacity=1, font=\fontsize{12}{18}\selectfont\sffamily, align=right] at (9.7, 3) {\textbf{\href{https://losdeldgiim.github.io/}{Los Del DGIIM}}};
            
            \node[anchor=south east, text=white, opacity=1, font=\fontsize{12}{15}\selectfont\sffamily, align=right] at (9.7, 1.8) {Doble Grado en Ingeniería Informática y Matemáticas\\Universidad de Granada};
        \end{tikzpicture}
    \end{figure}
    
    
    \restoregeometry        % Restaurar márgenes normales para las páginas subsiguientes
    \pagecolor{white}       % Restaurar el color de página
    
    
    \newpage
    \thispagestyle{empty}               % Sin encabezado ni pie de página
    \begin{tikzpicture}[remember picture, overlay]
        \node[anchor=south west, inner sep=3cm] at (current page.south west) {
            \begin{minipage}{0.5\paperwidth}
                \href{https://creativecommons.org/licenses/by-nc-nd/4.0/}{
                    \includegraphics[height=2cm]{assets/Licencia.png}
                }\vspace{1cm}\\
                Esta obra está bajo una
                \href{https://creativecommons.org/licenses/by-nc-nd/4.0/}{
                    Licencia Creative Commons Atribución-NoComercial-SinDerivadas 4.0 Internacional (CC BY-NC-ND 4.0).
                }\\
    
                Eres libre de compartir y redistribuir el contenido de esta obra en cualquier medio o formato, siempre y cuando des el crédito adecuado a los autores originales y no persigas fines comerciales. 
            \end{minipage}
        };
    \end{tikzpicture}
    
    
    
    % 1. Foto de fondo
    % 2. Título
    % 3. Encabezado Izquierdo
    % 4. Color de fondo
    % 5. Coord x del titulo
    % 6. Coord y del titulo
    % 7. Fecha


}


\newcommand{\portadaBook}[7]{

    % 1. Foto de fondo
    % 2. Título
    % 3. Encabezado Izquierdo
    % 4. Color de fondo
    % 5. Coord x del titulo
    % 6. Coord y del titulo
    % 7. Fecha

    % Personaliza el formato del título
    \pretitle{\begin{center}\bfseries\fontsize{42}{56}\selectfont}
    \posttitle{\par\end{center}\vspace{2em}}
    
    % Personaliza el formato del autor
    \preauthor{\begin{center}\Large}
    \postauthor{\par\end{center}\vfill}
    
    % Personaliza el formato de la fecha
    \predate{\begin{center}\huge}
    \postdate{\par\end{center}\vspace{2em}}
    
    \title{#2}
    \author{\href{https://losdeldgiim.github.io/}{Los Del DGIIM}}
    \date{Granada, #7}
    \maketitle
    
    \tableofcontents
}




\newcommand{\portadaArticle}[7]{

    % 1. Foto de fondo
    % 2. Título
    % 3. Encabezado Izquierdo
    % 4. Color de fondo
    % 5. Coord x del titulo
    % 6. Coord y del titulo
    % 7. Fecha

    % Personaliza el formato del título
    \pretitle{\begin{center}\bfseries\fontsize{42}{56}\selectfont}
    \posttitle{\par\end{center}\vspace{2em}}
    
    % Personaliza el formato del autor
    \preauthor{\begin{center}\Large}
    \postauthor{\par\end{center}\vspace{3em}}
    
    % Personaliza el formato de la fecha
    \predate{\begin{center}\huge}
    \postdate{\par\end{center}\vspace{5em}}
    
    \title{#2}
    \author{\href{https://losdeldgiim.github.io/}{Los Del DGIIM}}
    \date{Granada, #7}
    \thispagestyle{empty}               % Sin encabezado ni pie de página
    \maketitle
    \vfill
}
    \portadaExamen{ffccA4.jpg}{Ecuaciones\\Diferenciales I\\Examen IV}{Ecuaciones Diferenciales I. Examen IV}{MidnightBlue}{-8}{28}{2024-2025}{Arturo Olivares Martos}

    \begin{description}
        \item[Asignatura] Ecuaciones Diferenciales I
        \item[Curso Académico] 2016-17.
        %\item[Grado] Doble Grado en Ingeniería Informática y Matemáticas.
        \item[Grupo] B.
        \item[Profesor] Rafael Ortega Ríos.
        \item[Descripción] Primer parcial.
        \item[Fecha] 16 de marzo de 2017.
        %\item[Duración] 60 minutos.
    
    \end{description}
    \newpage

    \begin{ejercicio}
        Se considera la familia de curvas
        \begin{equation*}
            xy = c
        \end{equation*}
        donde $c \in \bb{R}$ es un parámetro. Calcula la familia de trayectorias ortogonales y dibuja las dos familias de curvas sobre el plano $(x, y)$.\\

        En primer lugar, necesitamos encontrar la ecuación diferencial que define a las curvas dadas. Derivando implícitamente, obtenemos:
        \begin{equation*}
            y + x\cdot y' = 0 \implies y' = -\frac{y}{x} \quad \text{con dominio } D=\left\{\begin{array}{c}
                D_1=\bb{R}^+\times\bb{R}\\ 
                \lor\\
                D_2=\bb{R}^-\times\bb{R}
            \end{array} \right.
        \end{equation*}

        Como sabemos que el producto de pendientes de rectas tangentes es $-1$, y sabiendo la interpretación geométrica de la derivada, podemos afirmar que las trayectorias ortogonales a las curvas dadas son las soluciones de la ecuación diferencial:
        \begin{equation*}
            y' = \frac{x}{y} \quad \text{con dominio } D=\left\{\begin{array}{rcl}
                D_{11} &=& \bb{R}^+\times\bb{R}^+\\
                D_{12} &=& \bb{R}^+\times\bb{R}^-\\
                D_{21} &=& \bb{R}^-\times\bb{R}^+\\
                D_{22} &=& \bb{R}^-\times\bb{R}^-
            \end{array} \right.
        \end{equation*}

        Resolvamos ahora esta ecuación diferencial. Se trata de una ecuación de variables separadas donde la función dependiente de $y$ no tiene soluciones constantes, por lo que:
        \begin{align*}
            \int y\,dy &= \int x\,dx\\
            \frac{y^2}{2} &= \frac{x^2}{2} + C'\\
            y^2 -x^2&= C
        \end{align*}

        Esto vemos que define una curva en implícitas $y(x)$ en el dominio correspondiente. Vemos que la familia de curvas dadas son las hipérbolas con ejes en los ejes coordenados, mientras que su curva de familias ortogonales son las hipérbolas con ejes en las rectas $y=\pm x$. Así, las dos familias de curvas dibujadas en el plano $(x, y)$ son:
        \begin{figure}[H]
            \centering
            \begin{tikzpicture}
                \begin{axis}[
                    axis lines = center,
                    xlabel = $x$,
                    ylabel = $y$,
                    xmin = -5,
                    ymin = -5,
                    xmax = 5,
                    ymax = 5,
                    xtick = \empty,
                    ytick = \empty
                ]
                % Dibujamos varias hipérbolas con ejes en los ejes coordenados, en distintos tonos de azul
                \addplot[domain=-4.5:4.5, samples=100, color=blue]{1/x};
                \addplot[domain=-4.5:4.5, samples=100, color=blue!70]{2/x};
                \addplot[domain=-4.5:4.5, samples=100, color=blue!25]{-1/x};

                % Dibujamos las hipérbolas con ejes en las rectas y=x, y=-x
                \addplot[domain=-4.5:4.5, samples=100, color=red]{sqrt(x^2+1)};
                \addplot[domain=-4.5:4.5, samples=100, color=red]{-sqrt(x^2+1)};

                \addplot[domain=-4.5:4.5, samples=100, color=red!50]{sqrt(x^2+4)};
                \addplot[domain=-4.5:4.5, samples=100, color=red!50]{-sqrt(x^2+4)};
                \end{axis}
            \end{tikzpicture}
        \end{figure}
    \end{ejercicio}

    \begin{ejercicio}
        Demuestra que la ecuación
        \begin{equation*}
            e^{x+t} + x = 0
        \end{equation*}
        define de forma implícita una función $x = x(t)$ que es derivable y está definida en todo $\bb{R}$. Encuentra una ecuación diferencial que admita a esta función como solución.\\

        Fijado $t\in \bb{R}$, definimos la función auxiliar
        \Func{f_t}{\bb{R}}{\bb{R}}{x}{e^{x+t} + x}

        Ver que esta ecuación define una función implícita $x(t)$ equivale a que la función $f_t$ tiene un único cero.
        \begin{description}
            \item[Existencia] En primer lugar, vemos que $f_t$ es continua. Además, tenemos que:
            \begin{equation*}
                \lim_{x\to-\infty} f_t(x) = -\infty \quad \text{y} \quad \lim_{x\to+\infty} f_t(x) = +\infty
            \end{equation*}

            Por tanto, por el teorema de Bolzano, existe $x(t)\in \bb{R}$ tal que $f_t(x(t)) = 0$.

            \item[Unicidad] Tenemos que $f_t\in C^1(\bb{R})$, con:
            \begin{equation*}
                f_t'(x) = e^{x+t} + 1 > 0 \quad \forall x\in\bb{R}
            \end{equation*}

            Por tanto, $f_t$ es estrictamente creciente, por lo que tiene a lo sumo un cero.
        \end{description}

        Por tanto, la ecuación dada define una función implícita $x(t)$ en todo $\bb{R}$:
        \Func{x}{\bb{R}}{\bb{R}}{t}{x(t)}

        Para ver que es derivable, hemos de aplicar el Teorema de la Función Implícita. Definimos la siguiente función auxiliar:
        \Func{F}{\bb{R}^2}{\bb{R}}{(t,x)}{e^{x+t} + x}

        Esta tiene derivadas parciales:
        \begin{equation*}
            \del{F}{t}(t,x) = e^{x+t} \quad \text{y} \quad \del{F}{x}(t,x) = e^{x+t} + 1
        \end{equation*}
        
        Ambas son continuas, por lo que $F\in C^1(\bb{R}^2)$. Además, la derivada parcial respecto de $x$ es siempre positiva (en particular, no se anula). Por tanto, para cada $t\in \bb{R}$, aplicamos el Teorema de la Función Implícita al punto $(t,x(t))$ (ya que $F(t,x(t)) = 0$). Obtenemos así que $x(t)\in C^1(\bb{R})$.\\

        Para encontrar una ecuación diferencial que admita a $x(t)$ como solución, derivamos implícitamente la ecuación dada:
        \begin{equation*}
            e^{x+t} + \left(e^{x+t} + 1\right)x' = 0 \implies x' = -\frac{e^{x+t}}{e^{x+t} + 1} \qquad \text{con dominio } D=\bb{R}^2
        \end{equation*}
    \end{ejercicio}


    \begin{ejercicio}
        Consideramos el siguiente sistema:
        \begin{equation*}
            \begin{cases}
                x' = y,\\
                y' = x
            \end{cases}
        \end{equation*}
        Este admite la solución $(x(t), y(t))$ con $x(t) = y(t) = e^t$, definida para todo $t \in \bb{R}$. Dibuja la órbita asociada en el plano $(x, y)$. Encuentra la ecuación diferencial de las órbitas de este sistema.\\

        En primer lugar, hemos de representar la órbita pedida. Como $x(t) = y(t)$, tenemos que la órbita está contenida en la recta $y = x$.
        No obstante, como tenemos que $x(\bb{R})=y(\bb{R}) = \left]0, +\infty\right[$, la órbita es la parte de la recta $y = x$ contenida en el primer cuadrante, sin incluir el origen. Esta es:
        \begin{figure}[H]
            \centering
            \begin{tikzpicture}
                \begin{axis}[
                    axis lines = center,
                    xlabel = $x$,
                    ylabel = $y$,
                    xmin = -1,
                    ymin = -1,
                    xmax = 5,
                    ymax = 5,
                    xtick = \empty,
                    ytick = \empty
                ]
                \addplot[domain=0:4.5, samples=2, color=red]{x};
                \node at (axis cs:0,0) [circle, draw, inner sep=3pt, color=red, ultra thick] {};
                \end{axis}
            \end{tikzpicture}
        \end{figure}

        Para encontrar la ecuación diferencial de las órbitas, hemos de considerar la derivada de la función $y = y(x)$ que define la órbita. Así, tenemos que:
        \begin{equation*}
            \frac{dy}{dx} = \frac{x}{y} \quad \text{con dominio } D=\left\{\begin{array}{rcl}
                D_1 &=& \bb{R}\times \bb{R}^+\\
                D_2 &=& \bb{R}\times \bb{R}^-
            \end{array} \right.
        \end{equation*}
        
    \end{ejercicio}

    \begin{ejercicio}
        Se considera el cambio de variables
        \begin{equation*}
            \varphi : s = x, y = t.
        \end{equation*}
        ¿En qué circunstancias se puede asegurar que es un cambio admisible para la ecuación diferencial $\frac{dx}{dt} = f(t, x)$ con $f : \bb{R}^2 \to \bb{R}$ continua?
        Se considera la nueva ecuación $\frac{dy}{ds} = \hat{f}(s, y)$. ¿Qué relación hay entre $f$ y $\hat{f}$?\\

        Para que el cambio de variables sea admisible, hemos de asegurar en primer lugar que es un difeomorfismo de clase $C^1$.
        Sea el cambio de variable $\varphi=(\varphi_1, \varphi_2)$:
        \Func{\varphi=(\varphi_1,\varphi_2)}{\bb{R}^2}{\bb{R}^2}{(t,x)}{(s,y)=(x,t)}
        Vemos que $\varphi$ se trata de la simetría axial respecto de la recta $y=x$. Comprobemos entonces que $\varphi$ es un difeomorfismo, para lo cual calculamos su inversa:
        \Func{\varphi^{-1}}{\bb{R}^2}{\bb{R}^2}{(s,y)}{(t,x)=(y,s)}

        Comprobemos que son inversas:
        \begin{align*}
            \varphi^{-1}(\varphi(t,x)) &= \varphi^{-1}(x,t) = (t,x)\qquad \forall (t,x)\in\bb{R}^2\\
            \varphi(\varphi^{-1}(s,y)) &= \varphi(y,s) = (s,y) \qquad \forall (s,y)\in\bb{R}^2
        \end{align*}

        Por tanto, $\varphi^{-1}$ es la inversa de $\varphi$, por lo que $\varphi$ es biyectiva. Además, como sus componentes son proyecciones, tenemos que $\varphi, \varphi^{-1}\in C^1(\bb{R}^2)$, por lo que $\varphi$ es un difeomorfismo.\\

        Calculemos ahora la condición de admisibilidad:
        \begin{align*}
            \del{\varphi_1}{t}(t,x) + \del{\varphi_1}{x}(t,x)f(t,x) &= 0 + 1\cdot f(t,x) = f(t,x) \neq 0 \qquad \forall (t,x)\in\bb{R}^2
        \end{align*}

        Por tanto, el cambio de variables es admisible si $f(t,x)\neq 0$ para todo $(t,x)\in\bb{R}^2$.\\

        Calculamos ahora la ecuación transformada. Tenemos que:
        \begin{equation*}
            \frac{dy}{ds} = \frac{dy}{dt}\cdot\frac{dt}{ds} = \frac{dt}{ds} = \frac{1}{x'} = \dfrac{1}{f(t,x)} = \dfrac{1}{f(y,s)} = \hat{f}(s,y)
            \qquad \text{con dominio } D=\bb{R}^2
        \end{equation*}

        Por tanto, tenemos que la relación entre $f$ y $\hat{f}$ es:
        \begin{equation*}
            \hat{f}(s,y) = \frac{1}{f(y,s)} \qquad \forall (s,y)\in\bb{R}^2
        \end{equation*}
    \end{ejercicio}

    \begin{ejercicio}
        Dada una función $F : \bb{R}^2 \to \bb{R}$, $(x, y) \mapsto F(x, y)$ de clase $C^1$ y un punto $(x_0, y_0) \in \bb{R}^2$ tal que $F(x_0, y_0) = 0$ se considera el problema de funciones implícitas
        \begin{equation*}
            F(x, y(x)) = 0, \quad y(x_0) = y_0.
        \end{equation*}
        Encuentra (si es que existen) una función $F$ y un punto $(x_0, y_0)$ en las condiciones anteriores y para los que el problema de funciones implícitas no admita solución.
        \begin{observacion}
            Se considera el problema local, la posible solución $y(x)$ está definida en algún entorno de $x_0$.
        \end{observacion}

        Consideramos la función $F(x, y) = x^2+y^2$, que es de clase $C^1(\bb{R}^2)$ por ser polinómica. Tomamos el punto $(x_0, y_0) = (0, 0)$, que cumple que $F(0, 0) = 0$.\\

        Para demostrar que no existe función implícita que resuelva el problema, el estudiante podría pensar en primer lugar en aplicar el Teorema de la Función Implícita y ver que no cumple las condiciones:
        \begin{equation*}
            \del{F}{y}(x_0, y_0) = 2y_0=0
        \end{equation*}
        No obstante, esto tan solo afirma que no podemos aplicar dicho Teorema. Recordemos que la implicación en el Teorema de la Función Implícita es de una sola dirección. Para comprobar que no existe función implícita, basta que demostrar que el único punto que cumple $F(x,y) = 0$ es el origen, por lo que no puede existir una función $y(x)$ que cumpla $F(x, y(x)) = 0$ en un entorno de $(0, 0)$.\\
    \end{ejercicio}
\end{document}